% =====================================================================
% Chapter 5 — Logic and Puzzles — part b: problems 14–30 (new problems, core level)
% =====================================================================
% a weight labelled with its mass (Problem 22)
\providecommand{\lgaWeight}[1]{\tikz[baseline=(lgaw.south)]{\node[trapezium,trapezium angle=72,draw=Slate,fill=Slate!25,line width=0.6pt,
  rounded corners=1.5pt,inner xsep=1pt,inner ysep=2.5pt,font=\scriptsize\bfseries] (lgaw) {#1};}}

% ---------------------------------------------------------------- Problem 14
\needspace{12\baselineskip}\begin{probbox}{14}{\relax}
On an island, every person either always tells the truth or always lies. You meet two people, Alice and Ben. Alice says:
``We are both liars.'' Which of the following is correct?
\choicesL{Both tell the truth.}{Both are liars.}{Alice tells the truth, Ben lies.}{Alice lies, Ben tells the truth.}{It cannot be determined.}
\end{probbox}
\probsol{
\textbf{Analysing the problem.} There is only one statement, Alice's. We consider two cases, Alice truthful or Alice lying, and see which case
gives no contradiction.
\begin{recallbox}A truth-teller only says true things; a liar only says false things. The negation of ``both are liars'' is
``at least one of them tells the truth''.\end{recallbox}
\step{1}{Case: Alice tells the truth}
Then Alice's statement is true: both are liars -- so Alice is a liar too. But we assumed Alice tells the truth. \textbf{Contradiction}, ruled out.
\step{2}{Case: Alice lies}
Then Alice's statement is false: it is \emph{not} the case that both are liars, so at least one of them tells the truth. Alice lies, so the
truth-teller must be \textbf{Ben}. No contradiction.
\step{3}{Conclusion}
Alice lies, Ben tells the truth.
\begin{proofbox}This is a version of ``I am a liar'': a truth-teller never calls himself a liar. So as soon as we hear Alice say ``we are both
liars'' (which includes Alice), we know Alice must be a liar.\end{proofbox}
\begin{warnbox}Many students choose (B), thinking ``Alice lies, so the opposite is true: both tell the truth''. Wrong! The opposite of ``both
are liars'' is only ``at least one tells the truth''.\end{warnbox}
\verify{Alice lies, Ben is truthful: the statement ``we are both liars'' is false (because Ben tells the truth) -- indeed a lie, which fits Alice.}
So the answer is \textbf{(D)}.
\ans{(D) Alice lies, Ben tells the truth}
}

% ---------------------------------------------------------------- Problem 15
\needspace{12\baselineskip}\begin{probbox}{15}{\relax}
A vase has been broken. Exactly one of the four friends David, Helen, Leo and Mia broke it. When asked, they answer:
\begin{itemize}[leftmargin=1.6em,itemsep=0pt,topsep=2pt]
\item David: ``Helen broke it.'' \qquad Helen: ``Leo broke it.''
\item Leo: ``Helen is lying.'' \qquad Mia: ``I did not break it.''
\end{itemize}
We know that \textbf{exactly one} of them is telling the truth. Who broke the vase?
\choices{David}{Helen}{Leo}{Mia}{Cannot tell}
\end{probbox}
\probsol{
\textbf{Analysing the problem.} Four statements, exactly one true, three false. We must find who broke the vase.
\begin{recallbox}For two contradictory statements (they cannot both be true and cannot both be false), exactly one of them is true.\end{recallbox}
\step{1}{Find a contradictory pair}
Leo says ``Helen is lying''. If Helen's statement is true, Leo's is false; if Helen's statement is false, Leo's is true. So among the statements of
Helen and Leo, \textbf{exactly one} is true.
\step{2}{So David and Mia are both lying}
The whole group has only one true statement, and it lies in the pair Helen -- Leo. Therefore the statements of David and Mia are both \textbf{false}.
\step{3}{Use Mia's false statement}
Mia's ``I did not break it'' is false, so \textbf{Mia broke the vase}.
\step{4}{Check everything}
\begin{center}
\begin{tabular}{|l|l|c|}
\hline
\rowcolor{tmBlue!10} Friend & Statement & True / false (if Mia broke it) \\ \hline
David & ``Helen broke it'' & false \\ \hline
Helen & ``Leo broke it'' & false \\ \hline
Leo & ``Helen is lying'' & \textbf{true} (Helen indeed said something false) \\ \hline
Mia & ``I did not'' & false \\ \hline
\end{tabular}
\end{center}
Exactly one true statement (Leo's). This matches the problem.
\begin{warnbox}Do not test people at random and stop as soon as something ``looks fine''. The sure way is to test all $4$ possibilities: if David
broke it, Leo's and Mia's statements are both true ($2$ true statements); if Helen broke it, David's, Leo's and Mia's are all true; if Leo broke it, Helen's
and Mia's are true. Only ``Mia broke it'' gives exactly one true statement.\end{warnbox}
So the answer is \textbf{(D)}.
\ans{(D) Mia}
}

% ---------------------------------------------------------------- Problem 16
\needspace{12\baselineskip}\begin{probbox}{16}{\relax}
Kate encrypts a short wish with a Caesar cipher with shift $5$ (each letter is replaced by the letter $5$ places after it on the circle
of letters A--Z, where A follows Z). She gets: \textbf{MFAJ KZS}. The original wish is
\choicesII{HAVE FUN}{RKFO PEX}{IBWF GVO}{HAVE FUM}{GZUD ETM}
\end{probbox}
\probsol{
\textbf{Analysing the problem.} We have the ciphertext and know the shift is $5$. To find the original we must \emph{decrypt}: shift
\emph{back} $5$ places.
\begin{recallbox}Decrypt: take the letter's number (A $= 0$, \ldots, Z $= 25$) and subtract $5$; if the result is negative, add $26$.\end{recallbox}
\step{1}{Decrypt each letter}
\begin{center}
\setlength{\tabcolsep}{4pt}
\begin{tabular}{|c|c|c|c|c|c|c|c|}
\hline
\rowcolor{tmBlue!10} Cipher letter & M & F & A & J & K & Z & S \\ \hline
Number & $12$ & $5$ & $0$ & $9$ & $10$ & $25$ & $18$ \\ \hline
Subtract $5$ & $7$ & $0$ & $-5 \to 21$ & $4$ & $5$ & $20$ & $13$ \\ \hline
Original & \textbf{H} & \textbf{A} & \textbf{V} & \textbf{E} & \textbf{F} & \textbf{U} & \textbf{N} \\ \hline
\end{tabular}
\end{center}
The letter A has to ``wrap around'': A $\to$ Z $\to$ Y $\to$ X $\to$ W $\to$ V (going back $5$ steps), which is exactly $-5 + 26 = 21 =$ V.
\step{2}{Read the result}
The wish is HAVE FUN.
\begin{warnbox}(B) RKFO PEX is what you get by shifting \emph{forwards} another $5$ -- encrypting once more instead of decrypting. (C) and (E)
are shifts back by $4$ and by $6$ (miscounting). (D) is wrong in the last letter: S shifted back $5$ is N, not M.\end{warnbox}
\verify{Encrypt HAVE FUN with shift $5$: H $\to$ M, A $\to$ F, V$(21) + 5 = 26 - 26 = 0 =$ A, E $\to$ J, F $\to$ K, U $\to$ Z, N $\to$ S.
We get MFAJ KZS.}
So the answer is \textbf{(A)}.
\ans{(A) HAVE FUN}
}

% ---------------------------------------------------------------- Problem 17
\needspace{12\baselineskip}\begin{probbox}{17}{\relax}
Nora encrypts the letter M with a Caesar cipher three times in a row: first with shift $9$, then with shift $11$ (applied to the result of
the first step), and finally with shift $9$ again. After the three steps, which letter has M become?
\choices{V}{G}{J}{P}{O}
\end{probbox}
\probsol{
\textbf{Analysing the problem.} Three shifts in the same direction. We can do them one at a time, or combine them into a single shift.
\begin{recallbox}Shifting by $a$ and then by $b$ is the same as shifting by $a + b$. A shift of $26$ returns to the start, so only the
remainder of the number of steps on division by $26$ matters.\end{recallbox}
\step{1}{Combine the three shifts}
Total steps: $9 + 11 + 9 = 29$. Since $29 = 26 + 3$, a shift of $29$ is exactly the same as a shift of $3$.
\step{2}{Shift M by $3$}
M $= 12$; $12 + 3 = 15 =$ P. (M $\to$ N $\to$ O $\to$ P.)
\step{3}{Check by doing one shift at a time}
\begin{center}
\begin{tabular}{|c|c|c|c|}
\hline
\rowcolor{tmBlue!10} Start & After step 1 ($+9$) & After step 2 ($+11$) & After step 3 ($+9$) \\ \hline
M $(12)$ & $21 =$ V & $32 - 26 = 6 =$ G & $15 =$ P \\ \hline
\end{tabular}
\end{center}
\begin{warnbox}Options (A) V and (B) G are \emph{intermediate} results (after steps $1$ and $2$). Read carefully: the question asks for the result
after \emph{all three} steps. (C) J is a shift \emph{backwards} by $3$.\end{warnbox}
\verify{Going back $29$ steps from P, i.e.\ back $3$ steps (since going back $26$ returns to the same letter): P $\to$ O $\to$ N $\to$ M. Correct.}
So the answer is \textbf{(D)}.
\ans{(D) P}
}

% ---------------------------------------------------------------- Problem 18
\needspace{12\baselineskip}\begin{probbox}{18}{\relax}
Five friends, Nina, Olivia, Peter, Quinn and Sam, all have different heights. Nina is taller than Olivia. Peter is shorter than Quinn but
taller than Nina. Sam is shorter than Olivia. When the five friends line up by height, who stands \emph{exactly in the middle}?
\choices{Nina}{Olivia}{Peter}{Quinn}{Sam}
\end{probbox}
\probsol{
\textbf{Analysing the problem.} There are five friends; ``in the middle'' is the third tallest (two taller and two shorter).
\step{1}{Use symbols (``$>$'' means ``taller than'')}
$N > O$; \quad $Q > P > N$; \quad $O > S$.
\step{2}{Chain them}
$Q > P > N$ and $N > O > S$ give: \quad $Q > P > N > O > S$.
\begin{figC}
  \draw[->,Navy,line width=0.8pt] (0,-0.3)--(0,2.8) node[above,font=\scriptsize,text=Slate]{taller};
  \foreach \y/\n in {2.4/Quinn,1.8/Peter,1.2/Nina,0.6/Olivia,0/Sam}{
    \fill[tmBlue] (0,\y) circle (2.2pt); \node[right,font=\small] at (0.15,\y){\n};}
  \node[anchor=west,font=\small\bfseries,text=AnsRed] at (1.45,1.2){$\leftarrow$ in the middle (3rd)};
\end{figC}
\step{3}{Read off the result}
From tallest to shortest: Quinn, Peter, \textbf{Nina}, Olivia, Sam. The person in the middle is Nina.
\begin{warnbox}``Peter is shorter than Quinn but taller than Nina'' contains \emph{two} facts. Many students write only $Q > P$ and forget
$P > N$; then the chain cannot be completed.\end{warnbox}
\verify{Nina (3rd) is taller than Olivia (4th) \checkmark; Peter (2nd) is shorter than Quinn and taller than Nina \checkmark; Sam (5th) is shorter than Olivia \checkmark.}
So the answer is \textbf{(A)}.
\ans{(A) Nina}
}

% ---------------------------------------------------------------- Problem 19
\needspace{12\baselineskip}\begin{probbox}{19}{\relax}
In a queue for tickets, Henry is $7$th from the front and $12$th from the back. How many people are in the queue?
\choices{$17$}{$18$}{$19$}{$20$}{$5$}
\end{probbox}
\probsol{
\textbf{Analysing the problem.} We know Henry's position counted from both ends. We must find the total number of people in the queue.
\step{1}{Draw a diagram}
\begin{figC}
  \foreach \i in {1,...,18}{\fill[tmBlue!60] (\i*0.5,0) circle (0.15);}
  \fill[AnsRed] (3.5,0) circle (0.17);
  \node[font=\scriptsize,text=AnsRed] at (3.5,0.4){Henry};
  \draw[decorate,decoration={brace,amplitude=4pt,mirror},tmGreen!60!black,line width=0.8pt] (0.35,-0.3)--(3.65,-0.3) node[midway,below=4pt,font=\scriptsize]{$7$ people (including Henry)};
  \draw[decorate,decoration={brace,amplitude=4pt},tmOrange!80!black,line width=0.8pt] (3.35,0.75)--(9.15,0.75) node[midway,above=4pt,font=\scriptsize]{$12$ people (including Henry)};
  \node[font=\scriptsize,text=Slate] at (0.5,0.4){front}; \node[font=\scriptsize,text=Slate] at (9,0.4){back};
\end{figC}
\step{2}{Count}
Counting from the front to Henry: $7$ people. Counting from Henry to the back: $12$ people. Adding gives $7 + 12 = 19$, but Henry has been counted
\textbf{twice}. So the queue has $19 - 1 = 18$ people.
\begin{recallbox}Being $k$-th from the front and $m$-th from the back means the line has $k + m - 1$ people.\end{recallbox}
\step{3}{Another way}
In front of Henry there are $7 - 1 = 6$ people, behind him $12 - 1 = 11$ people. The whole queue: $6 + 1 + 11 = 18$ people.
\begin{warnbox}Option (C) $19 = 7 + 12$ counts Henry twice; option (E) $5 = 12 - 7$ takes the difference -- which means nothing here.\end{warnbox}
\verify{In a queue of $18$ people, if Henry is $7$th from the front, then from the back he is $18 - 7 + 1 = 12$th. Correct.}
So the answer is \textbf{(B)}.
\ans{(B) $18$}
}

% ---------------------------------------------------------------- Problem 20
\needspace{12\baselineskip}\begin{probbox}{20}{\relax}
\begin{minipage}[t]{0.63\linewidth}
Place the numbers $1, 2, 3, \ldots, 9$ (each once) in the grid shown to make a magic square (the numbers in every row, every column and both
diagonals have the same sum). Three numbers are already filled in. Complete the grid. What is the number in the bottom-left corner?
\end{minipage}\hfill
\begin{minipage}[t]{0.33\linewidth}\vspace{-4pt}\centering
\lgaMagic{$2$}{$7$}{}{}{}{$1$}{}{}{}
\end{minipage}
\end{probbox}
\probsol{
\textbf{Analysing the problem.} A $3\times 3$ magic square with the numbers $1$ to $9$; we know the top-left cell ($2$), the middle of the
top row ($7$) and the right cell of the middle row ($1$). We must fill the $6$ remaining cells.
\step{1}{Magic sum and centre}
\begin{recallbox}$S = (\text{sum of the } 9 \text{ numbers}) \div 3$; centre $= S \div 3$.\end{recallbox}
$1 + 2 + \cdots + 9 = 45$, so $S = 45 \div 3 = 15$; centre $= 15 \div 3 = 5$.
\step{2}{Fill in the lines that already have two numbers}
\begin{itemize}[leftmargin=1.4em,itemsep=1pt,topsep=1pt]
\item Top row: $2 + 7 + \square = 15$ $\Rightarrow$ top-right corner $= 6$.
\item Middle row: $\square + 5 + 1 = 15$ $\Rightarrow$ left cell of the middle row $= 9$.
\item Diagonal $2$ -- $5$ -- ?: bottom-right corner $= 15 - 2 - 5 = 8$.
\item Diagonal $6$ -- $5$ -- ?: bottom-left corner $= 15 - 6 - 5 = \mathbf{4}$.
\item Middle column: $7 + 5 + \square = 15$ $\Rightarrow$ middle of the bottom row $= 3$.
\end{itemize}
\begin{center}\lgaMagic{$2$}{$7$}{\textcolor{tmGreen!60!black}{$6$}}{\textcolor{tmGreen!60!black}{$9$}}{\textcolor{tmGreen!60!black}{$5$}}{$1$}{\textcolor{AnsRed}{$\mathbf{4}$}}{\textcolor{tmGreen!60!black}{$3$}}{\textcolor{tmGreen!60!black}{$8$}}\end{center}
\verify{Rows: $2 + 7 + 6 = 15$, $9 + 5 + 1 = 15$, $4 + 3 + 8 = 15$. Columns: $2 + 9 + 4 = 15$, $7 + 5 + 3 = 15$, $6 + 1 + 8 = 15$.
Diagonals: $2 + 5 + 8 = 15$, $6 + 5 + 4 = 15$. The numbers $1$ to $9$ are each used exactly once.}
\begin{warnbox}You must check \emph{all} eight lines (three rows, three columns, two diagonals) and make sure no number is repeated. A grid
that works for rows and columns but not for a diagonal is not a magic square.\end{warnbox}
\ans{Grid: $2, 7, 6$ / $9, 5, 1$ / $4, 3, 8$. The bottom-left number is $4$.}
}

% ---------------------------------------------------------------- Problem 21
\needspace{12\baselineskip}\begin{probbox}{21}{\relax}
\begin{minipage}[t]{0.63\linewidth}
The numbers $3, 4, 5, \ldots, 11$ (nine consecutive whole numbers) are placed in the grid shown, each once, to make a magic square.
The number $x$ in the bottom-right corner is
\choicesII{$6$}{$8$}{$10$}{$3$}{$5$}
\end{minipage}\hfill
\begin{minipage}[t]{0.33\linewidth}\vspace{-4pt}\centering
\lgaMagic{$4$}{$9$}{}{$11$}{}{}{}{}{$x$}
\end{minipage}
\end{probbox}
\probsol{
\textbf{Analysing the problem.} Nine numbers from $3$ to $11$. Three cells are known; the cell $x$ lies on the main diagonal together with
the cell $4$ and the centre.
\step{1}{Find the magic sum}
Sum of the nine numbers: $3 + 4 + \cdots + 11$. Pair first with last: $(3 + 11) + (4 + 10) + (5 + 9) + (6 + 8) + 7 = 4\times 14 + 7 = 63$.
So $S = 63 \div 3 = 21$.
\step{2}{The centre}
\begin{recallbox}In a $3\times 3$ magic square, centre $= S \div 3$.\end{recallbox}
Centre $= 21 \div 3 = 7$ (which is also the middle number of the list $3, 4, \ldots, 11$).
\step{3}{Find $x$ on the main diagonal}
$4 + 7 + x = 21$ $\Rightarrow$ $x = 21 - 11 = 10$.
\step{4}{Complete the grid to check}
Top row: $21 - 4 - 9 = 8$. Middle row: $21 - 11 - 7 = 3$. Left column: $21 - 4 - 11 = 6$. Bottom row: $21 - 6 - 10 = 5$.
\begin{center}\lgaMagic{$4$}{$9$}{\textcolor{tmGreen!60!black}{$8$}}{$11$}{\textcolor{tmGreen!60!black}{$7$}}{\textcolor{tmGreen!60!black}{$3$}}{\textcolor{tmGreen!60!black}{$6$}}{\textcolor{tmGreen!60!black}{$5$}}{\textcolor{AnsRed}{$\mathbf{10}$}}\end{center}
\verify{Middle column $9 + 7 + 5 = 21$; right column $8 + 3 + 10 = 21$; other diagonal $8 + 7 + 6 = 21$. All nine numbers $3$ to $11$ are used, no repeats.}
\begin{warnbox}If you mistakenly take the magic sum to be $15$ (as for the numbers $1$ to $9$), you get $x = 4$ -- a number that is already
used! The magic sum depends on which numbers are used.\end{warnbox}
So the answer is \textbf{(C)}.
\ans{(C) $10$}
}

% ---------------------------------------------------------------- Problem 22
\needspace{12\baselineskip}\begin{probbox}{22}{\relax}
Both scales below are balanced. The circles all have the same mass.
\begin{center}
\lgaScale{\lgaCir\,\lgaCir\,\lgaCir}{\lgaSq}\hspace{1cm}\lgaScale{\lgaSq\,\lgaCir}{\lgaWeight{800 g}}
\end{center}
How many grams does the square \lgaSq\ weigh?
\end{probbox}
\probsol{
\textbf{Analysing the problem.} There are two unknown quantities: the mass of one circle and the mass of the square. The two scales give
two equalities.
\begin{recallbox}If a group of objects $A$ balances a group $B$, then on another scale we may replace $A$ by $B$ (substitution).\end{recallbox}
\step{1}{Choose the unknown}
Let the mass of one circle be $x$ grams, $x > 0$. The first scale tells us that the square weighs as much as three circles, so the square
weighs $3x$ grams.
\step{2}{Form an equation from the second scale}
The left pan holds the square and one circle: $3x + x$ grams; the right pan holds the $800$ g weight. So
\[ 3x + x = 800, \qquad\text{that is,}\qquad 4x = 800. \]
\step{3}{Solve the equation}
$x = 800 \div 4 = 200$. One circle weighs $200$ g.
\step{4}{Answer the question}
The square weighs $3x = 3\cdot 200 = 600$ g.
\begin{warnbox}The question asks for the mass of the \emph{square}, not the circle. Stopping at $x = 200$ answers the wrong question.\end{warnbox}
\verify{Scale 1: $3\cdot 200 = 600$ g, the mass of the square \checkmark. Scale 2: $600 + 200 = 800$ g \checkmark.}
\ans{The square weighs $600$ g.}
}

% ---------------------------------------------------------------- Problem 23
\needspace{12\baselineskip}\begin{probbox}{23}{\relax}
There are four kinds of boxes: red \lgaBox{AnsRed}, yellow \lgaBox{StarGold}, blue \lgaBox{tmBlue}, purple \lgaBox{tmPurple}; boxes of the same
colour have equal masses. The following three scales are all balanced:
\begin{center}
\lgaScale{\lgaBox{AnsRed}\,\lgaBox{StarGold}}{\lgaBox{tmBlue}}\hspace{0.4cm}\lgaScale{\lgaBox{tmBlue}\,\lgaBox{StarGold}}{\lgaBox{tmPurple}}\hspace{0.4cm}\lgaScale{\lgaBox{AnsRed}\,\lgaBox{AnsRed}}{\lgaBox{tmBlue}}
\end{center}
How many yellow boxes does one purple box weigh?
\choices{$2$}{$3$}{$4$}{$5$}{$6$}
\end{probbox}
\probsol{
\textbf{Analysing the problem.} Let the masses of the red, yellow, blue and purple boxes be $R, Y, B, P$. The three scales give three equalities:
\[ R + Y = B, \qquad B + Y = P, \qquad R + R = B. \]
We must express $P$ in terms of $Y$.
\step{1}{Compare scales 1 and 3}
Both have a single blue box on the right pan, so their left pans weigh the same: $R + Y = R + R$.
Remove one red box from each side: $Y = R$. \emph{A red box weighs the same as a yellow box.}
\step{2}{The blue box}
$B = R + Y = Y + Y = 2Y$: a blue box weighs as much as $2$ yellow boxes.
\step{3}{The purple box}
$P = B + Y = 2Y + Y = 3Y$: a purple box weighs as much as $3$ yellow boxes.
\begin{figC}
  \node at (0,0){\lgaScale{\lgaBox{tmPurple}}{\lgaBox{StarGold}\,\lgaBox{StarGold}\,\lgaBox{StarGold}}};
\end{figC}
\begin{warnbox}Do not assume a red box is heavier than a yellow one just because ``red'' appears more often. Use the scales to \emph{prove}
the relationships.\end{warnbox}
\verify{Take $Y = R = 1$ kg, $B = 2$ kg, $P = 3$ kg: scale 1: $1 + 1 = 2$; scale 2: $2 + 1 = 3$; scale 3: $1 + 1 = 2$. All three balance.}
So the answer is \textbf{(B)}.
\ans{(B) $3$}
}

% ---------------------------------------------------------------- Problem 24
\needspace{12\baselineskip}\begin{probbox}{24}{\relax}
\begin{minipage}[t]{0.63\linewidth}
The numbers $1, 2, 3, 4, 5, 6, 7, 8$ are placed in the eight circles of the square shown (four at the corners, four in the middles of the
sides), each once, so that the three numbers on each side add up to $13$. The sum of the four corner numbers is
\choices{$12$}{$13$}{$16$}{$20$}{$36$}
\end{minipage}\hfill
\begin{minipage}[t]{0.33\linewidth}\vspace{0pt}\centering
\lgaSqEight{}{}{}{}{}{}{}{}
\end{minipage}
\end{probbox}
\probsol{
\textbf{Analysing the problem.} The square has $4$ sides, each made of two corners and one middle number. We must find the sum of the four corners.
\begin{recallbox}When we add the sums of all the sides, a corner number is counted twice (it belongs to two sides) and a middle number once.\end{recallbox}
\step{1}{Add the sums of the four sides}
Each side sums to $13$, so the four sides together give $4\times 13 = 52$.
\step{2}{Count how many times each number is counted}
In those $52$, each number from $1$ to $8$ is counted at least once, and the four corner numbers are counted once more:
\[ 52 = (1 + 2 + \cdots + 8) + (\text{sum of the four corners}) = 36 + (\text{sum of the four corners}). \]
\step{3}{Compute}
Sum of the four corners $= 52 - 36 = 16$.
\step{4}{An actual arrangement (to be sure the puzzle has a solution)}
\begin{figC}
  \node at (0,0){\lgaSqEight{$1$}{$8$}{$4$}{$3$}{$6$}{$2$}{$5$}{$7$}};
  \node[font=\small,anchor=west,align=left] at (2.0,0){top: $1 + 8 + 4 = 13$\\ right: $4 + 3 + 6 = 13$\\ bottom: $6 + 2 + 5 = 13$\\ left: $5 + 7 + 1 = 13$\\[2pt] corners: $1 + 4 + 6 + 5 = 16$};
\end{figC}
\begin{warnbox}Option (E) $36$ is the sum of \emph{all} eight numbers; option (B) $13$ is the sum of one side. The question asks for the sum of the
\emph{corner} numbers.\end{warnbox}
So the answer is \textbf{(C)}.
\ans{(C) $16$}
}

% ---------------------------------------------------------------- Problem 25
\needspace{12\baselineskip}\begin{probbox}{25}{\relax}
\begin{minipage}[t]{0.63\linewidth}
Place the numbers $1, 2, 3, 4, 5$ in the five cells of the plus-shaped figure shown, each once, so that the sum of the three numbers in the
horizontal row equals the sum of the three numbers in the vertical column. How many different values can the centre cell take?
\choices{$1$}{$2$}{$3$}{$4$}{$5$}
\end{minipage}\hfill
\begin{minipage}[t]{0.33\linewidth}\vspace{-6pt}\centering
\lgaCross{}{}{$?$}{}{}
\end{minipage}
\end{probbox}
\probsol{
\textbf{Analysing the problem.} The centre cell (call it $c$) belongs to both the row and the column. The other four cells form two pairs: the
left -- right pair (row) and the top -- bottom pair (column).
\step{1}{The condition for equal sums}
Row: $(\text{left} + \text{right}) + c$; column: $(\text{top} + \text{bottom}) + c$. The two sums are equal exactly when
\[ \text{left} + \text{right} = \text{top} + \text{bottom}. \]
\step{2}{The sum of the other four numbers must be even}
The four non-centre numbers add up to $15 - c$ (since $1 + 2 + 3 + 4 + 5 = 15$). The two pairs have equal sums, so $15 - c$ must split into two
equal halves, i.e.\ $15 - c$ is even. Since $15$ is odd, $c$ must be \textbf{odd}: $c \in \{1, 3, 5\}$.
\step{3}{Check each odd value}
\begin{itemize}[leftmargin=1.4em,itemsep=1pt,topsep=1pt]
\item $c = 1$: $2, 3, 4, 5$ remain, each pair must sum to $14 \div 2 = 7$: \ $2 + 5 = 3 + 4 = 7$. \checkmark
\item $c = 3$: $1, 2, 4, 5$ remain, each pair must sum to $12 \div 2 = 6$: \ $1 + 5 = 2 + 4 = 6$. \checkmark
\item $c = 5$: $1, 2, 3, 4$ remain, each pair must sum to $10 \div 2 = 5$: \ $1 + 4 = 2 + 3 = 5$. \checkmark
\end{itemize}
\begin{figC}
  \node at (0,0){\lgaCross{$3$}{$2$}{$1$}{$5$}{$4$}}; \node[font=\small] at (0,-1.5){$c = 1$};
  \node at (3.2,0){\lgaCross{$2$}{$1$}{$3$}{$5$}{$4$}}; \node[font=\small] at (3.2,-1.5){$c = 3$};
  \node at (6.4,0){\lgaCross{$2$}{$1$}{$5$}{$4$}{$3$}}; \node[font=\small] at (6.4,-1.5){$c = 5$};
\end{figC}
\step{4}{Why does an even $c$ not work?}
If $c = 2$: the remaining numbers $1, 3, 4, 5$ add up to $13$ -- odd, so they cannot form two pairs with equal sums. If $c = 4$: $1, 2, 3, 5$
remain, with sum $11$ -- odd, also impossible.
\begin{warnbox}Do not think ``the centre must be the middle number ($3$)''. Here there are three possibilities!\end{warnbox}
So the centre can be $1$, $3$ or $5$: $3$ values. The answer is \textbf{(C)}.
\ans{(C) $3$}
}

% ---------------------------------------------------------------- Problem 26
\needspace{12\baselineskip}\begin{probbox}{26}{\relax}
\begin{minipage}[t]{0.66\linewidth}
Three standard dice (opposite faces add up to $7$) are stacked on a table as shown. The top face shows $2$ pips. The total number of pips on
the five hidden \emph{horizontal} faces (the faces pressed together between two dice and the bottom face resting on the table) is
\choices{$14$}{$17$}{$19$}{$21$}{$23$}
\end{minipage}\hfill
\begin{minipage}[t]{0.3\linewidth}\vspace{-4pt}\centering
\begin{tikzpicture}[scale=0.62,line join=round]
  \draw[Slate,line width=1pt] (-0.4,0)--(2.0,0);
  \foreach \k/\f/\r in {0/4/1,1/5/3,2/3/1}{
    \begin{scope}[yshift=\k cm]
    \draw[Ink,line width=0.8pt,fill=white] (0,0) rectangle (1,1);
    \draw[Ink,line width=0.8pt,fill=white] (1,0)--(1.45,0.45)--(1.45,1.45)--(1,1)--cycle;
    \begin{scope}[Ink]\lgaPips{\f}\end{scope}
    \begin{scope}[Ink,cm={0.45,0.45,0,1,(1,0)}]\lgaPips{\r}\end{scope}
    \end{scope}}
  \draw[Ink,line width=0.8pt,fill=white] (0,3)--(1,3)--(1.45,3.45)--(0.45,3.45)--cycle;
  \begin{scope}[Ink,cm={1,0,0.45,0.45,(0,3)}]\lgaPips{2}\end{scope}
\end{tikzpicture}
\end{minipage}
\end{probbox}
\probsol{
\textbf{Analysing the problem.} The five hidden horizontal faces are: the bottom face of the top die; the top and bottom faces of the middle
die; the top and bottom faces of the bottom die. We do not need to know each face separately -- we only use ``opposite faces add up to $7$''.
\begin{recallbox}On a standard die the top and bottom faces are opposite, so together they carry $7$ pips.\end{recallbox}
\step{1}{Bottom die and middle die}
For each of them the top and bottom faces are both hidden, and each pair totals $7$. Two dice: $7 + 7 = 14$.
\step{2}{Top die}
The visible top face has $2$ pips, so the (opposite) bottom face has $7 - 2 = 5$ pips, and it is hidden.
\step{3}{Add up}
$14 + 5 = 19$.
\begin{warnbox}(D) $21 = 3\times 7$ forgets that the very top face ($2$ pips) \emph{can be seen}; it is not hidden. (A) $14$ forgets the bottom
face of the top die (also hidden, with $5$ pips).\end{warnbox}
\verify{The six horizontal faces of the three dice total $3\times 7 = 21$; of these only the very top face ($2$ pips) is visible, so the five hidden ones total $21 - 2 = 19$.}
So the answer is \textbf{(C)}.
\ans{(C) $19$}
}

% ---------------------------------------------------------------- Problem 27
\needspace{12\baselineskip}\begin{probbox}{27}{\relax}
\begin{minipage}[t]{0.63\linewidth}
The $4\times 4$ grid shown is to be filled so that every row, every column and every thick-bordered $2\times 2$ box contains the digits
$1, 2, 3, 4$. The digit in cell $Q$ is
\choicesL{$1$}{$2$}{$3$}{$4$}{It cannot be determined}
\end{minipage}\hfill
\begin{minipage}[t]{0.33\linewidth}\vspace{4pt}\centering
\lgaFour{2/1/$3$,4/1/$2$,4/2/$1$,1/3/$4$,3/4/$2$,4/4/$Q$}
\end{minipage}
\end{probbox}
\probsol{
\textbf{Analysing the problem.} Cell $Q$ is in row $4$, column $4$ (bottom-right corner), in the bottom-right $2\times 2$ box.
\begin{recallbox}The digit in a cell must differ from every digit already in its row, its column and its $2\times 2$ box.\end{recallbox}
\step{1}{Column $4$}
Column $4$ already has $2$ (row $1$) and $1$ (row $2$). The two remaining cells -- row $3$ column $4$ and cell $Q$ -- must be $3$ and $4$ in
some order.
\step{2}{Row $3$ decides}
Row $3$ already has a $4$ (in column $1$), so row $3$ column $4$ \emph{cannot} be $4$. Hence row $3$ column $4$ is $3$, and $Q = 4$.
\step{3}{Check with the bottom-right box and row $4$}
The bottom-right box contains $3$ (row $3$ column $4$), $2$ (row $4$ column $3$), $Q = 4$, and the remaining cell (row $3$ column $3$) is
$1$ -- all of $1, 2, 3, 4$. Row $4$ has $2$ and $4$, no repeat.
\step{4}{Complete the grid (to be sure there is no contradiction)}
\begin{figC}
  \node at (0,0){\lgaFour{1/1/$1$,2/1/$3$,3/1/$4$,4/1/$2$,1/2/$2$,2/2/$4$,3/2/$3$,4/2/$1$,1/3/$4$,2/3/$2$,3/3/$1$,4/3/$3$,1/4/$3$,2/4/$1$,3/4/$2$,4/4/\textcolor{AnsRed}{$\mathbf{4}$}}};
  \node[font=\small,anchor=west,align=left] at (1.8,0){Row $1$: $1\,3\,4\,2$\\ Row $2$: $2\,4\,3\,1$\\ Row $3$: $4\,2\,1\,3$\\ Row $4$: $3\,1\,2\,4$};
\end{figC}
\verify{Every row, column and $2\times 2$ box contains all of $1, 2, 3, 4$; a computer check confirms this is the only way to fill the grid.}
\begin{warnbox}Looking only at column $4$, $Q$ could be $3$ or $4$, and it is tempting to choose (E). Also look at row $3$ -- a different
region containing a ``neighbour'' of $Q$.\end{warnbox}
So the answer is \textbf{(D)}.
\ans{(D) $4$}
}

% ---------------------------------------------------------------- Problem 28
\needspace{12\baselineskip}\begin{probbox}{28}{\relax}
In class 7A, \emph{all} members of the chess club can swim. Leo cannot swim. Nora can swim. Which of the following statements \textbf{must} be true?
\choicesL{Leo is in the chess club.}{Leo is not in the chess club.}{Nora is in the chess club.}{Nora is not in the chess club.}{Everyone in the class who can swim is in the chess club.}
\end{probbox}
\probsol{
\textbf{Analysing the problem.} ``All members of the chess club can swim'' has the form ``every A is a B'' (A: in the chess club, B: can swim).
\begin{recallbox}From ``every A is a B'' we may conclude ``not a B means not an A''; but we may \emph{not} conclude ``a B means an A''.\end{recallbox}
\step{1}{Draw a diagram}
\begin{figC}
  \draw[tmBlue,line width=1pt,fill=tmBlue!8] (0.9,1.1) ellipse (1.9 and 1.2);
  \draw[tmOrange,line width=1pt,fill=tmOrange!15] (0.3,1.0) ellipse (0.95 and 0.7);
  \node[font=\scriptsize,text=tmOrange!70!black] at (0.3,1.3){chess club};
  \node[font=\scriptsize,text=tmBlue!80!black] at (1.95,1.6){can swim};
  \fill[tmGreen!60!black] (1.95,0.95) circle (2.2pt) node[below,font=\scriptsize]{Nora?};
  \fill[tmGreen!60!black] (0.3,0.9) circle (2.2pt) node[below,font=\scriptsize]{Nora?};
  \fill[AnsRed] (3.6,1.1) circle (2.2pt) node[right,font=\scriptsize]{Leo};
\end{figC}
The ``chess club'' oval lies entirely inside the ``can swim'' oval.
\step{2}{Leo}
Leo cannot swim, so Leo is \emph{outside} the big oval, and therefore certainly outside the small one: \textbf{Leo is not in the chess club}.
If Leo were in the club, Leo would have to be able to swim -- contradicting the problem (proof by contradiction).
\step{3}{Nora}
Nora can swim: Nora is inside the big oval, but may be inside or outside the small one. So (C) and (D) are each only \emph{possibly} true, not certain.
\step{4}{Statement (E)}
(E) reverses the given statement. The problem says nothing about pupils who can swim but do not play chess, so (E) is not certainly true.
\begin{warnbox}A very common mistake: from ``chess players can all swim'' concluding ``swimmers all play chess''. Think of an example: ``every cat
has four legs'' does not mean ``every four-legged animal is a cat''!\end{warnbox}
So the answer is \textbf{(B)}.
\ans{(B) Leo is not in the chess club}
}

% ---------------------------------------------------------------- Problem 29
\needspace{12\baselineskip}\begin{probbox}{29}{\relax}
Mr Fox lies on Mondays, Tuesdays and Wednesdays and tells the truth on the other days of the week. One day, Fox says: ``Yesterday I lied.''
Which day could it be?
\choicesL{Only Monday.}{Only Thursday.}{Monday or Thursday.}{Wednesday or Thursday.}{Sunday or Monday.}
\end{probbox}
\probsol{
\textbf{Analysing the problem.} Fox can say ``Yesterday I lied'' when:
\begin{itemize}[leftmargin=1.4em,itemsep=0pt,topsep=1pt]
\item he tells the truth today and the statement is true: truthful today, lied yesterday; or
\item he lies today and the statement is false: lying today, truthful yesterday.
\end{itemize}
That is: \textbf{yesterday and today are of different types}.
\begin{recallbox}The day before Monday is Sunday (the week goes round).\end{recallbox}
\step{1}{Make a table}
\begin{center}
\setlength{\tabcolsep}{4pt}
\begin{tabular}{|l|c|c|c|c|c|c|c|}
\hline
\rowcolor{tmBlue!10} Today & Mon & Tue & Wed & Thu & Fri & Sat & Sun \\ \hline
Fox today & L & L & L & T & T & T & T \\ \hline
Fox yesterday & T (Sun) & L & L & L & T & T & T \\ \hline
Different? & \textbf{yes} & -- & -- & \textbf{yes} & -- & -- & -- \\ \hline
\end{tabular}
\end{center}
\step{2}{Check the two days found}
\emph{Monday:} Fox lies; yesterday (Sunday) he told the truth, so ``yesterday I lied'' is false -- indeed a lie. \checkmark\\
\emph{Thursday:} Fox tells the truth; yesterday (Wednesday) he lied, so the statement is true -- indeed the truth. \checkmark
\begin{warnbox}Wednesday: Fox lies, and yesterday (Tuesday) he lied too, so ``yesterday I lied'' is a \emph{true} statement. Someone who is lying
does not say true things, so Wednesday is ruled out -- this is the trap in option (D).\end{warnbox}
So the answer is \textbf{(C)}.
\ans{(C) Monday or Thursday}
}

% ---------------------------------------------------------------- Problem 30
\needspace{12\baselineskip}\begin{probbox}{30}{\relax}
Four friends, Amy, Ben, Chloe and David, run a race; no two finish at the same time. Three spectators make predictions:
\begin{itemize}[leftmargin=1.6em,itemsep=0pt,topsep=2pt]
\item Henry: ``Amy comes first, Ben comes second.''
\item Kate: ``Chloe comes first, David comes fourth.''
\item Lucy: ``Ben comes third, David comes second.''
\end{itemize}
In the end, each spectator got \emph{exactly one} part right. Find the finishing order of the four friends.
\end{probbox}
\probsol{
\textbf{Analysing the problem.} Each prediction has two parts, exactly one of them right. We split into cases using Henry's prediction.
\begin{recallbox}``Half-true'' statement: if one part is true the other is false, and if one part is false the other is true.\end{recallbox}
\step{1}{Case 1: ``Ben second'' is true (and ``Amy first'' is false)}
Look at Lucy's prediction. Ben is second, so ``Ben third'' is false; hence the other part ``David second'' must be true. But Ben is already
second, so David cannot be second too. \textbf{Contradiction} -- this case is ruled out.
\step{2}{Case 2: ``Amy first'' is true (and ``Ben second'' is false)}
\begin{itemize}[leftmargin=1.4em,itemsep=1pt,topsep=1pt]
\item Kate: Amy is first, so ``Chloe first'' is false $\Rightarrow$ ``David fourth'' is true.
\item Lucy: David is fourth, so ``David second'' is false $\Rightarrow$ ``Ben third'' is true.
\item Second place is left for Chloe.
\end{itemize}
Order: Amy (1st), Chloe (2nd), Ben (3rd), David (4th).
\step{3}{Check}
\begin{center}
\begin{tabular}{|l|c|c|}
\hline
\rowcolor{tmBlue!10} Spectator & First part & Second part \\ \hline
Henry & Amy first: \textcolor{tmGreen!70!black}{true} & Ben second: \textcolor{AnsRed}{false} \\ \hline
Kate & Chloe first: \textcolor{AnsRed}{false} & David fourth: \textcolor{tmGreen!70!black}{true} \\ \hline
Lucy & Ben third: \textcolor{tmGreen!70!black}{true} & David second: \textcolor{AnsRed}{false} \\ \hline
\end{tabular}
\end{center}
\begin{tipbox}Choose the prediction to split on so that each case ``triggers'' as many other parts as possible. Here ``Ben second'' clashes
immediately with Lucy's ``David second'', so case 1 is ruled out very quickly.\end{tipbox}
\begin{parentbox}You can use four slips of paper with the names, line them up as ``1st -- 2nd -- 3rd -- 4th'' and swap them around to test.
Then ask your child: ``Have we tried every case? How can we be sure there is no other answer?'' -- the answer is Step 1.\end{parentbox}
\ans{Amy 1st, Chloe 2nd, Ben 3rd, David 4th.}
}
